%==============================================================
% Tensor Tree Sketching
%
% Compile with:  pdflatex -> bibtex -> pdflatex -> pdflatex
% (or simply:    latexmk -pdf main.tex)
%==============================================================
\documentclass[11pt]{article}

%--------------------------------------------------------------
% Encoding and fonts
%--------------------------------------------------------------
\usepackage[T1]{fontenc}
\usepackage[utf8]{inputenc}
\usepackage{xcolor}

%--------------------------------------------------------------
% Math
%--------------------------------------------------------------
\usepackage{amsmath}
\usepackage{amssymb}
\usepackage{amsthm}
\usepackage{mathtools}
\usepackage{bm}

%--------------------------------------------------------------
% Graphics and misc
%--------------------------------------------------------------
\usepackage{graphicx}
\usepackage{tikz}
\usetikzlibrary{positioning, calc}
\usepackage{xcolor}
\usepackage{subcaption}
\usepackage{algorithm,algorithmic}

%--------------------------------------------------------------
% Paper margins
%--------------------------------------------------------------
\usepackage{geometry}
\geometry{verbose,tmargin=3.5cm,bmargin=3.5cm,lmargin=2.5cm,rmargin=2.5cm}

%--------------------------------------------------------------
% References (biblatex with the classic bibtex backend, matching refs.bib)
%--------------------------------------------------------------
\usepackage[
  backend=bibtex,
  style=numeric,
  sorting=nyt,
  maxbibnames=99,
  giveninits=true,
]{biblatex}
\addbibresource{refs.bib}

%--------------------------------------------------------------
% Hyperlinks and cross-references
% (hyperref then cleveref, in this order)
%--------------------------------------------------------------
\usepackage{hyperref}
\usepackage[capitalize]{cleveref}
\crefname{definition}{Definition}{Definitions}
\Crefname{definition}{Definition}{Definitions}

%--------------------------------------------------------------
% Project macros
%--------------------------------------------------------------
\newcommand{\R}{\mathbb{R}}
\newcommand{\C}{\mathbb{C}}
\newcommand{\F}{\mathbb{F}}
\newcommand{\Z}{\mathbb{Z}}
\newcommand{\N}{\mathbb{N}}

\newcommand{\si}{\mathrm{i}}
\newcommand{\rpart}{\mathop{\mathrm{Re}}}
\newcommand{\ipart}{\mathop{\mathrm{Im}}}

% Top operators on matrices and scalars
\newcommand{\T}{\textsc{T}}
\newcommand{\pinv}{+}
\newcommand{\winv}{-}
\newcommand{\ha}{\textsc{H}}
\newcommand{\sconj}{*}
\newcommand{\adj}{*}

% Other matrix products
\newcommand{\kronprod}{\otimes}
\newcommand{\hadprod}{\bullet}

% Norms and inner products
\newcommand{\inner}[2]{\left\langle #1,#2 \right\rangle}
\newcommand{\dotprod}[2]{#1\cdot #2}
\newcommand{\Xnorm}[2]{\left\| #1 \right\|_{#2}}
\newcommand{\Fnorm}[1]{\Xnorm{#1}{F}}
\newcommand{\Tnorm}[1]{\Xnorm{#1}{2}}
\newcommand{\XnormS}[2]{\left\| #1 \right\|^2_{#2}}
\newcommand{\FnormS}[1]{\XnormS{#1}{F}}
\newcommand{\TnormS}[1]{\XnormS{#1}{2}}

\newcommand{\reshape}{\mathop{\mathrm{reshape}}}
\newcommand{\rank}{\mathop{\mathrm{rank}}}
\newcommand{\diag}{\mathop{\mathrm{diag}}}
\newcommand{\trace}{\mathop{\mathrm{Tr}}}
\newcommand{\vectorize}{\mathop{\mathrm{vec}}}

% Modifiers for vector, matrix and tensor
\renewcommand{\vec}[1]{\mathbf{\boldsymbol{#1}}}
\newcommand{\mat}[1]{\mathbf{#1}}
\newcommand{\ten}[1]{\mathcal{#1}}

% Vector symbols
\newcommand{\va}{\vec{a}}
\newcommand{\vb}{\vec{b}}
\newcommand{\vc}{\vec{c}}
\newcommand{\vd}{\vec{d}}
\newcommand{\ve}{\vec{e}}
\newcommand{\vi}{\vec{i}}
\newcommand{\vr}{\vec{r}}
\newcommand{\vs}{\vec{s}}
\newcommand{\vt}{\vec{t}}
\newcommand{\vq}{\vec{q}}
\newcommand{\vu}{\vec{u}}
\newcommand{\vv}{\vec{v}}
\newcommand{\vw}{\vec{w}}
\newcommand{\vx}{\vec{x}}
\newcommand{\vy}{\vec{y}}
\newcommand{\vz}{\vec{z}}

\newcommand{\vzero}{\vec{0}}
\newcommand{\valpha}{\vec{\alpha}}
\newcommand{\vbeta}{\vec{\beta}}
\newcommand{\vgamma}{\vec{\gamma}}
\newcommand{\vtheta}{\vec{\theta}}
\newcommand{\veta}{\vec{\eta}}
\newcommand{\vxi}{\vec{\xi}}
\newcommand{\vsigma}{\vec{\sigma}}
\newcommand{\vlambda}{\vec{\lambda}}
\newcommand{\vrho}{\vec{\rho}}
\newcommand{\vpsi}{\vec{\psi}}
\newcommand{\vphi}{\vec{\phi}}

% Matrix symbols
\newcommand{\matA}{\mat{A}}
\newcommand{\matB}{\mat{B}}
\newcommand{\matC}{\mat{C}}
\newcommand{\matD}{\mat{D}}
\newcommand{\matE}{\mat{E}}
\newcommand{\matF}{\mat{F}}
\newcommand{\matG}{\mat{G}}
\newcommand{\matH}{\mat{H}}
\newcommand{\matK}{\mat{K}}
\newcommand{\matM}{\mat{M}}
\newcommand{\matN}{\mat{N}}
\newcommand{\matO}{\mat{O}}
\newcommand{\matP}{\mat{P}}
\newcommand{\matQ}{\mat{Q}}
\newcommand{\matR}{\mat{R}}
\newcommand{\matS}{\mat{S}}
\newcommand{\matT}{\mat{T}}
\newcommand{\matI}{\mat{I}}
\newcommand{\matJ}{\mat{J}}
\newcommand{\matU}{\mat{U}}
\newcommand{\matV}{\mat{V}}
\newcommand{\matW}{\mat{W}}
\newcommand{\matX}{\mat{X}}
\newcommand{\matY}{\mat{Y}}
\newcommand{\matZ}{\mat{Z}}

\newcommand{\matSigma}{\mat{\Sigma}}
\newcommand{\matPi}{\mat{\Pi}}
\newcommand{\matLambda}{\mat{\Lambda}}
\newcommand{\matTheta}{\mat{\Theta}}

\newcommand{\matZero}{\mat{0}}
\newcommand{\matOnes}{\mat{1}}

% Tensor symbols
\newcommand{\tenA}{\ten{A}}
\newcommand{\tenB}{\ten{B}}
\newcommand{\tenC}{\ten{C}}
\newcommand{\tenD}{\ten{D}}
\newcommand{\tenE}{\ten{E}}
\newcommand{\tenI}{\ten{I}}
\newcommand{\tenU}{\ten{U}}
\newcommand{\tenV}{\ten{V}}
\newcommand{\tenS}{\ten{S}}
\newcommand{\tenX}{\ten{X}}
\newcommand{\tenY}{\ten{Y}}
\newcommand{\tenZ}{\ten{Z}}
\newcommand{\tenT}{\ten{T}}

% General tensor algebra
\newcommand{\squeeze}{\mathop{\mathrm{squeeze}}}
\newcommand{\twist}{\mathop{\mathrm{twist}}}
\newcommand{\fold}{\mathop{\mathrm{fold}}}
\newcommand{\unfold}{\mathop{\mathrm{unfold}}}
\newcommand{\nmodeprod}[1]{\times_{#1}}
\newcommand{\vnmodeprod}[1]{\bar{\times}_{#1}}
\newcommand{\modesplit}[2]{#1_{[#2]}}
\newcommand{\contract}[2]{\mathop{\mathrm{Tr}_{#1,#2}}}
\newcommand{\faceprod}{\Delta}
\newcommand{\ctrprod}{\bowtie}
\newcommand{\outprod}{\circ}

% Tubal tensor algebra
\newcommand{\tprod}{\star_{\text{t}}}
\newcommand{\Mprod}{\star_{\matM}}
\newcommand{\Mleq}{\leq_{\matM}}
\newcommand{\Mgeq}{\geq_{\matM}}
\newcommand{\unitM}{\ve_{\matM}}
\newcommand{\KM}{{\mathbb{K}}_{\matM}}
\newcommand{\Mdot}{{\cdot}_{\matM}}

%--------------------------------------------------------------
% Lean-formalization annotations
% Visible while drafting; define \LEANPRODUCTION to hide them.
% See lean/CLAUDE.md for the paper<->Lean status workflow.
%--------------------------------------------------------------
\input{lean/latex/leanrefs}

%--------------------------------------------------------------
% Theorem-like environments
%--------------------------------------------------------------
\theoremstyle{plain}
\newtheorem{theorem}{Theorem}[section]
\newtheorem{lemma}[theorem]{Lemma}
\newtheorem{proposition}[theorem]{Proposition}
\newtheorem{corollary}[theorem]{Corollary}

\theoremstyle{definition}
\newtheorem{definition}[theorem]{Definition}
\newtheorem{assumption}[theorem]{Assumption}
\newtheorem{example}[theorem]{Example}

\theoremstyle{remark}
\newtheorem{remark}[theorem]{Remark}

%==============================================================
% Title
%==============================================================
\title{Tensor Tree Sketching}

\author{%
  Guy Ben Eliezer \and
  Leia Greenberg \and
  Haim Avron\thanks{Tel Aviv University}
  % \and Co-author \thanks{Affiliation.}
}

\date{\today}

%==============================================================
\begin{document}

\maketitle
\tableofcontents
\leanmodebanner
%--------------------------------------------------------------
% Introduction
%--------------------------------------------------------------

\section{Introduction}
\label{sec:introduction}

In many computational tasks the object of interest is the trace of a symmetric positive
semidefinite (PSD) matrix $\matA$ presented not as an explicit array but as
a \emph{tensor network}, a collection of smaller tensors and their connections whose contraction
represents $\matA$. A useful feature of this representation is that the trace $\trace(\matA)$ can be computed efficiently at the network level. We circumvent the explicit construction of $\matA$. This matters precisely when assembling $\matA$ is impossible {\color{red} [Add citations]}.
For instance, in many body quantum physics {\color{red} [Add citations]} a matrix product
operator (MPO) is a tensor network representing a matrix operating on a quantum system of $N$ sites represented in itself as another tensor network. We get that the matrix represented by an MPO for a system with $N$  sites where each site has a physical dimension $d$ is $\matA\in\C^{d^{N}\times d^{N}}$, i.e. is exponentially large in $N$ and cannot be stored entirely in memory. We are then forced to access $\matA$ only through algebraic operations as $\vx^{\top}\matA\vx$ which are computable efficiently at the network level. 
To approximate the trace, we can use a Girard-Hutchinson estimator (\cite{girard1989, hutchinson1989}): $$\trace(\matA)\approx\frac1\ell\sum_{i=1}^l\vx_{i}^{\top}\matA\vx_{i}$$ 
where the vectors $\vx_{i}, i\in[l]$  are independent, identically distributed random vectors with the isotropic property, i.e. $\mathbb{E}[\vx^T\vx]=\matI$.

A central instance arises in quantum many-body physics. We need to calculate the \emph{Reny moments} for a $N$ sized subsystem inside a larger general tensor network, $G$ with $mat(G)=\mat{A}\in \C^{d^N\times d^N}$, i.e.  we calculate $Tr(\matA^k)$. In some systems the calculation of $\vx^T\matA^k\vx$  has an exponential cost in $k$, which bring us back to the same point. 
There are methods {\color{red} [cite Noa]} that circumvent this by sampling $k$  Kronecker structure tensors $\vx^{(j)}=\vy_1^{(j)}\kronprod\cdots \vy_N^{(j)}$ where $\vx^{(j)}\in \C^{d^N},\vy_n^{(j)}\in \C^d$ with $j=1\in[k]$  and calculating $$\trace(\matA^k)\approx\frac1\ell\sum_{i=1}^l\vx_{i}^{(1)\top}\matA\vx_{i}^{(2)}\vx_{i}^{(2)\top}\matA\vx_i^{(3)}\cdots \vx_i^{(k)\top}\matA\vx_i^{(1)}$$
Those methods reduce $\trace(\matA^k)$  into $k$  scalars and multiply, circumventing the need for the exponentially large in $k$ tensors. With that said, they still have exponential sample complexity in $N$ and $k$ {\color{red} [cite Haim]}.

We seek an estimator that is simultaneously  efficient to contract and whose sample complexity is polynomial in all system parameters.
Instead of taking a tensor such that $\vx^T\matA\vx$  is a scalar, we randomly sample sketching tensor network, $G_{S^{(i)}}, i\in[k]$  with $mat(G_S)=\matS\in\R^{m\times d^N}$ where we call $m$  the \emph{sketch dimension}. We then sketch our data $\matA$ to matrices $\matS^{(i)T}\mathbf{AS}^{(i+1)}\in \R^{m\times m}$ and multiply the matrices, and we finally calculate $$\trace(\matA^k)\approx \trace\!\left(\matS^{(1)}\matA\,\matS^{(2)\top}
      \matS^{(2)}\matA\cdots
      \matS^{(k)}\matA\,\matS^{(1)\top}\right)$$   We present a theoretical framework that establishes the precise conditions under which the property holds for $(\varepsilon,\delta)$ accurate embedding: 
      \begin{equation}
          \Pr[|\trace(\matA^k)-\trace\!\left(\matS^{(1)}\matA\,\matS^{(2)\top}
      \matS^{(2)}\matA\cdots
      \matS^{(k)}\matA\,\matS^{(1)\top}\right)|>\varepsilon\cdot\trace(\matA^k)]<\delta
      \end{equation}
\paragraph{Our Contributions}
\begin{itemize}
  \item \textbf{Theory (general).} For any PSD matrix $\matA$ in a tensor-network form, we find a family of sketches that gives the estimator  \[
    E_{k}
      =\trace\!\left(\matS_{1}\matA\matS_{2}^{\top}
        \matS_{2}\matA\cdots\matS_{k}\matA\matS_{1}^{\top}\right)
  \]
  With polynomial sketch size in amount of network dangling edges, edge dimensions and power $k$ (Theorem \ref{thm:tracek}).
  \item \textbf{General Tensor Network Algorithm.} We give a heuristic algorithm to contract a general symmetric tensor network and calculate our estimator. We show that in the MPOP case we give better results by achieving a smaller error with a lower computational cost (see Sec. \ref{sec:alg-mpo})
  \item  \textbf{Numerical Results.} We report numerical experiments to assert our claim for our theorem and algorithms for the MPOP and general case. 
\end{itemize}

\subsection{Related Work}
\label{sec:related}

{\color{red} TODO} 
The threads to be developed are: 
\begin{itemize}
\item stochastic trace estimation and its
variance-reduced variants \cite{hutchinson1989,girard1989,meyer2021hutchpp,persson2022nystrompp}

\item Kronecker and rank-one probing, with its lower bounds
\cite{meyer2026kronecker,bujanovic2021rankone,feldman2022entanglement}

\item tensor-train random vectors for trace estimation \cite{bujanovic2026tt}

and structured tensor-network sketches
\cite{masolomonik2022,cazeaux2026linear,rakhshan2021rademacher,oseledets2011tt}.

\item A recurring point is that oblivious subspace embeddings
\cite{cazeaux2026linear,camano2025osi} are the wrong primitive here, since for a
full-rank $\matA$ the relevant subspace is all of $\R^{d^{N}}$
\item our element-wise union bound instead pays only the logarithm of the eigenbasis
size (\Cref{rem:osi}).

\end{itemize}

\section{Preliminaries}
\subsection{Notation}
We use lower case bold letters for vectors $({\mathbf{v,w}},...)$ upper case bold letters for matrices $(\mathbf{X,Y,..})$ and calligraphic letters for higher order tensors $(\tenA,\tenB)$. For $n\in\N$  we denote $[n]\coloneqq\{1,\dots n\}$. The identity matrix of size n is written as $\matI_{n}$. For matrices $\matX,\matY $, we denote $\matX\kronprod\matY$  as the Kronecker product. 

\subsection{Definitions}
{\color{blue} Leia's note: i think most (or all) of the defs should go in an appendix.}

\begin{definition}[Tensor]\label[definition]{def:tensor}
\leanref{TensorTreeSketching.Tensor}\leanok
Let $\mathbb{F}$ denote the field of scalars, either the real numbers
$\mathbb{R}$ or the complex numbers $\mathbb{C}$. An order-$n$ tensor
$\tenA\in\mathbb{F}^{s_1\times\cdots\times s_n}$ is an $n$-dimensional
array, with entries $\tenA_{i_1,\dots,i_n}\in\mathbb{F}$ for
$i_k\in[s_k]$.
\end{definition}

\begin{definition}[Contraction]\label[definition]{def:contraction}
\leanref{TensorTreeSketching.contract2}\leanok
Let $\tenA\in\mathbb{F}^{a_1\times\cdots\times a_p\times s_1\times\cdots\times s_t}$
and
$\tenB\in\mathbb{F}^{s_1\times\cdots\times s_t\times b_1\times\cdots\times b_q}$
share the $t$ modes of sizes $s_1,\dots,s_t$, indexed by the common multi-index
$\vr=(r_1,\dots,r_t)\in[s_1]\times\cdots\times[s_t]$. Their
\emph{contraction} $(\tenA,\tenB)$ is the order-$(p+q)$ tensor in
$\mathbb{F}^{a_1\times\cdots\times a_p\times b_1\times\cdots\times b_q}$ obtained
by summing the product over the shared index set,
\[
  (\tenA,\tenB)_{\vi,\vec{j}}
    =\sum_{\vr\in[s_1]\times\cdots\times[s_t]}
      \tenA_{\vi,\vr}\,\tenB_{\vr,\vec{j}},
  \qquad
  \vi\in[a_1]\times\cdots\times[a_p],\ \
  \vec{j}\in[b_1]\times\cdots\times[b_q],
\]
whose modes are the unshared modes of $\tenA$ followed by those of
$\tenB$.
\end{definition}

\begin{remark}[Computational Cost of a contraction]\label[remark]{rem:contraction-cost}
\leanref{TensorTreeSketching.contractCost}\leanok
With $\tenA,\tenB$ as in Definition~\ref{def:contraction}, unfolding
$\tenA$ to
$\mathbb{F}^{(\prod_{u=1}^{p}a_u)\times(\prod_{w=1}^{t}s_w)}$ and $\tenB$ to
$\mathbb{F}^{(\prod_{w=1}^{t}s_w)\times(\prod_{v=1}^{q}b_v)}$, the contraction is
the product of these two matrices, with computational cost
\[
  C_{\tenA,\tenB}
    =\Bigl(\prod_{u=1}^{p}a_u\Bigr)
     \Bigl(\prod_{w=1}^{t}s_w\Bigr)
     \Bigl(\prod_{v=1}^{q}b_v\Bigr).
\]
The cost of contracting an entire network is the sum of the pairwise costs
$C_{\tenA,\tenB}$ incurred along a chosen contraction order, depending on that order.
\end{remark}

\begin{definition}[Tensor network]\label[definition]{def:tn}
A \emph{tensor network} is an undirected graph (or hyper-graph) $G=(V,E)$ in which each vertex
$v\in V$ carries a tensor $\tenA_v$ whose modes are represented by the
edges incident to $v$. An edge shared by two vertices identifies two modes of
equal size, which are summed over. The tensor represented by $G$ is the
tensor $\tenA$ obtained by contracting every such shared edge; we write
$G=G_{\tenA}$.
\end{definition}

\begin{definition}[Dangling edge]\label[definition]{def:dangling}
A \emph{dangling edge} of $G$ is an edge with a single endpoint. Its mode is not
summed over and remains a mode of the represented tensor. Denoting the dangling
edges $\{e_1,\dots,e_n\}$, the order of $\tenA$ equals the number of
dangling edges, so $\tenA$ is an order-$n$ tensor whose modes are exactly
$e_1,\dots,e_n$.
\end{definition}

\begin{definition}[Intermediate tensor and its edges]\label[definition]{def:intermediate}
Let $G=(V,E)$ be a tensor network and $W\subseteq V$. The \emph{intermediate
tensor of $W$}, denoted $\tenA_W$, is the tensor represented by the sub
network of the vertices in $W$: every edge with both endpoints in $W$ is
contracted, while every dangling edge incident to a vertex of $W$ and every
edge joining $W$ to $V\setminus W$ remains a mode of $\tenA_W$. We write
$E(W)$ for this set of remaining edges (the modes of $\tenA_W$), and
$E(W,W')$ for the set of edges with one endpoint in $W$ and the other in a
disjoint set $W'\subseteq V$.
\end{definition}

\begin{definition}[Contraction path]\label[definition]{def:cpath}
A \emph{contraction path} of a tensor network $G=(V,E)$ is a sequence of
$|V|-1$ contractions
$T_0=\bigl((U_1,W_1),\dots,(U_{|V|-1},W_{|V|-1})\bigr)$
in which, starting from each vertex $v\in V$ as its own intermediate tensor,
each step $(U_i,W_i)$ contracts two distinct intermediate tensors, on the
vertex sets $U_i$ and $W_i$, into a single intermediate on $U_i\cup W_i$
(\Cref{def:contraction}), so the last step returns a single tensor.
\end{definition}

\begin{definition}[Contraction tree of a path]\label[definition]{def:ctree}
The \emph{contraction tree} $T$ of a contraction path $T_0$ of $G=(V,E)$ is
the rooted binary tree whose leaves are the singletons $\{v\}$, $v\in V$, and
in which, for every contraction $(U,W)$ of $T_0$, the vertex $U\cup W$ has
children $U$ and $W$. It is a full binary tree with $|V|-1$ internal vertices
and root $V$.
Note that every topological sort (children before parents) of a contraction
tree is a contraction path.
\end{definition}
\begin{definition}[Constrained contraction tree]\label[definition]{def:constrained}
Let $G_{\tenA}=(V_{\tenA},E_{\tenA})$ be a data tensor network with
contraction path $T_0$ (\Cref{def:cpath}) and contraction tree $T$
(\Cref{def:ctree}), and let $G=(V_{\tenA}\cup V_S,E_{\tenA}\cup E_S)$ be the
joint network obtained by attaching an embedding network $G_S=(V_S,E_S)$ to
dangling edges of $G_{\tenA}$ - e.g.\ samples of a sketching tensor network
(\Cref{def:sketching-tn}) attached to the row and column edges. A contraction
tree $T_B$ of $G$ is constrained on $T_0$, written $T_0\subset T_B$,
if for every contraction $(U,W)\in T$ there exists a contraction
$(\hat U,\hat W)\in T_B$ with
\[
  \hat U\cap V_{\tenA}=U
  \qquad\text{and}\qquad
  \hat W\cap V_{\tenA}=W
\]
\cite[Def.~1]{masolomonik2022}.
\end{definition}

\begin{remark}[Cost of a contraction path]\label[remark]{rem:path-cost}
Let $T_0=\{(U_1,W_1),\dots,(U_{|V|-1},W_{|V|-1})\}$ be a contraction path of
a tensor network $G_{\tenA}=(V,E)$. Each contraction $(U_i,W_i)$ is a
contraction of two tensors in the sense of Definition~\ref{def:contraction},
namely of the tensors represented by the sub networks on $U_i$ and $W_i$,
and carries the cost $C_{U_i,W_i}$ of Remark~\ref{rem:contraction-cost}. The
\emph{total cost} of the path is
\[
  C(T_0)\;=\;\sum_{i=1}^{|V|-1} C_{U_i,W_i},
\]
and depends on the choice of path.
\end{remark}

\begin{definition}[Symmetric tensor]\label[definition]{def:symtensor}
\leanref{TensorTreeSketching.IsSymmetricMat}\leanok
A \emph{symmetric tensor} is an order-$2N$ tensor $\tenA$ with modes
$(e_1^{r},\dots,e_N^{r})$ and $(e_1^{c},\dots,e_N^{c})$, representing the rows
and columns with $\dim(e_i^{r})=\dim(e_i^{c})$ for all $i\in[N]$, such that when
the modes are ordered as
$(e_1^{r},\dots,e_N^{r},e_1^{c},\dots,e_N^{c})$, we get
\[
  \tenA_{i_1,\dots,i_N,\,j_1,\dots,j_N}
    =\tenA_{j_1,\dots,j_N,\,i_1,\dots,i_N}
  \qquad\text{for all index values.}
\]
\end{definition}

\begin{remark}\label[remark]{rem:symorder}
This definition is invariant under a common permutation of the rows and columns:
reordering the pairs $(e_i^{r},e_i^{c})$ by the same $\sigma\in S_N$ in both
blocks leaves symmetry unchanged. We fix one listing of the pairs throughout.
\end{remark}

\begin{definition}[Dimension tree]\label[definition]{def:dimtree}
Let $T$ be a contraction tree of $G_{\tenA}$. The
\emph{row dimension tree} of $T$, denoted $T_r(T)$, is the rooted tree whose
vertices are subsets of $\bar E^{\,r}$, whose leaves are the singletons
$\{e_i^{r}\}$, $i\in[N]$, and whose root is $\bar E^{\,r}$, in which two
vertices $E_1',E_2'$ share the parent $E_1'\cup E_2'$ if and only if there
exists a contraction $(U,W)\in T$ with
\[
  \bar E^{\,r}(U)=E_1'
  \qquad\text{and}\qquad
  \bar E^{\,r}(W)=E_2'.
\]
The \emph{column dimension tree} $T_c(T)$ is defined symmetrically with
respect to $\bar E^{\,c}$.
\end{definition}

\begin{figure}[htb]
  \centering
  \includegraphics[width=\textwidth]{Illustrations/symmetric_tn.png}
  \caption{A symmetric tensor network and its trees.
  (a)~A tensor network $G_{\tenA}$ on six vertices with row edges
  $e_1^{r},e_2^{r},e_3^{r}$ and column edges $e_1^{c},e_3^{c},e_2^{c}$;
  the planar positions of $e_2^{c}$ and $e_3^{c}$ are exchanged, so the
  symmetry for \Cref{def:symtn} is preserved when looking at the actual contraction path.
  (b)~A contraction tree of $G_{\tenA}$ (\Cref{def:ctree}); the circled
  numbers give a contraction path, i.e.
  $T_0=\bigl((v_2,v_3),(\{v_2,v_3\},v_1),(u_2,u_3),(\{u_2,u_3\},u_1),
  (\{v_1,v_2,v_3\},\{u_1,u_2,u_3\})\bigr)$.
  (c)~The resulting dimension tree (\Cref{def:dimtree}): step~1 merges
  $\{e_2^{r}\},\{e_3^{r}\}$ and step~3 merges
  $\{e_3^{c}\},\{e_2^{c}\}=\{e_2^{c}\},\{e_3^{c}\}$, so the row and
  column dimension trees coincide under $e_i^{r}\leftrightarrow e_i^{c}$,
  exactly the requirement of \Cref{def:symtn}}
  \label{fig:symtn}
\end{figure}


\begin{definition}[Symmetric tensor network]\label[definition]{def:symtn}
We call the tuple ($G_\tenA$, $(e_i^{r})_{i\in[N]}$, $(e_i^{c})_{i\in[N]}$, $T_0$) where $G_{\tenA}=(V_\tenA,E_\tenA)$ is a tensor network represent a symmetric tensor with dangling edges $(e_i^{r})_{i\in[N]}$, $(e_i^{c})_{i\in[N]}$, and $T_0$ be a contraction path, and let $T$ be the dimension tree of $(G_{\tenA},T_0)$. We
call  ($G_\tenA$, $(e_i^{r})_{i\in[N]}$, $(e_i^{c})_{i\in[N]}$, $T_0$) a \emph{symmetric tensor network} if the row and column dimension trees are equal up to exchanging for all, in the sense that each node $\{e_{i_1}^{r},\dots,e_{i_l}^{r}\}$ of the row dimension tree has a node in the column dimension tree $\{e_{i_1}^{c},\dots,e_{i_l}^{c}\}$ is a node and vice versa. See fig (\ref{fig:symtn}).
\end{definition}

\begin{figure}[H]
  \centering
  \includegraphics[width=0.5\textwidth]{Illustrations/sketching_tn.png}
  \caption{A sketching tensor network $G_S$ (gray, \Cref{def:sketching-tn})
  with $N=2$ dangling edges, attached to the row edges of the network of
  \Cref{fig:symtn}(a). Its structure follows the dimension tree of
  \Cref{fig:symtn}(c): $\tenS_{t_1}$ merges $e_2^{r}$ and $e_3^{r}$,
  and $\tenS_{t_2}$ merges the result with $e_1^{r}$ and carries the
  sketch dimension $m$ (red).}
  \label{fig:sketchtn}
\end{figure}


\begin{definition}[Sketching tensor network]\label[definition]{def:sketching-tn}
Let $\bigl(G_{\tenA},(e_i^{r})_{i\in[N]},(e_i^{c})_{i\in[N]},T_0\bigr)$
be a symmetric tensor network with $d_i=\dim(e_i^{r})$.
A \emph{sketching tensor network} for it is a tuple
\[
  \bigl(G_S=(V_S,E_S),\;\mu,\;m,\;(e_i^{s})_{i\in[N]}\bigr),
\]
where $G_S$ is a tensor network whose tensor entries are i.i.d.\ random
variables drawn from the distribution $\mu$, and whose dangling edges are $e_1^{s},\dots,e_N^{s}$ with one additional dangling edge of dimension $m$, called the \emph{sketch dimension}, with $\dim(e_i^{s})=d_i$ for all $i\in[N]$.
The matricization $\matS=\mathrm{mat}(G_S)\in\mathbb{F}^{m\times \prod_{i=1}^N d_i}$ groups the sketch dimension into the row index and $e_1^{s},\dots,e_N^{s}$, in order, into the column index. 
A \emph{sample} of the sketching tensor network is a realization of the entries of its tensors, drawn i.i.d.\ from $\mu$. Attaching a sample to the row edges of $G_{\tenA}$ means
identifying $e_i^{s}$ with $e_i^{r}$ for all $i\in[N]$, and analogously
for the column edges.
\end{definition}

\begin{definition}[Constrained contraction tree]\label[definition]{def:constrained}
Let $G_{\tenA}=(V_{\tenA},E_{\tenA})$ be a data tensor network with
contraction path $T_0$ (\Cref{def:cpath}) and contraction tree $T$
(\Cref{def:ctree}), and let $G=(V_{\tenA}\cup V_S,E_{\tenA}\cup E_S)$ be the
joint network obtained by attaching an embedding network $G_S=(V_S,E_S)$ to
dangling edges of $G_{\tenA}$, i.e. samples of a sketching tensor network
(\Cref{def:sketching-tn}) attached to the row and column edges. A contraction
tree $T_B$ of $G$ is \emph{constrained on} $T_0$, written $T_0\subset T_B$,
if for every contraction $(U,W)\in T$ there exists a contraction
$(\hat U,\hat W)\in T_B$ with
\[
  \hat U\cap V_{\tenA}=U
  \qquad\text{and}\qquad
  \hat W\cap V_{\tenA}=W
\]
\cite[Def.~1]{masolomonik2022}.
\end{definition}

\begin{definition}[Trace of a tensor network]\label[definition]{def:trace}
\leanref{TensorTreeSketching.trace_matricization}\leanok
Let $\tenA$ be an order-$2N$ tensor represented by a network
$G_{\tenA}$ with dangling edges paired as
$(e_i^{r})_{i\in[N]}$, $(e_i^{c})_{i\in[N]}$ and $\dim(e_i^{r})=\dim(e_i^{c})$
for all $i$. The \emph{trace} of $\tenA$ is the scalar obtained by
identifying $e_i^{r}$ with $e_i^{c}$ for every $i\in[N]$ and contracting the
resulting network, which has no dangling edges:
\[
  \trace(\tenA)
    =\sum_{i_1,\dots,i_N}\tenA_{i_1,\dots,i_N,\,i_1,\dots,i_N}.
\]
\end{definition}

\begin{definition}[Power of a symmetric tensor network]\label[definition]{def:power}
% The closing-the-power claim (Tr of the closed k-th power = Tr(A^k)) is what the
% Lean corollary formalizes; the former standalone cor:trace-power folded here.
\leanref{TensorTreeSketching.trace_power_eq_sum_diag}\leanok
Fix an integer $k\ge2$, and let $\tenA^{(1)},\dots,\tenA^{(k)}$ be
$k$ copies of a symmetric tensor $\tenA$, the copy $\tenA^{(\ell)}$
carrying dangling edges $(e_i^{r,(\ell)})_{i\in[N]}$,
$(e_i^{c,(\ell)})_{i\in[N]}$. The \emph{$k$-th power} of $\tenA$ is the
order-$2N$ tensor obtained by identifying, for every $i\in[N]$, the column edges
of each copy with the row edges of the next,
\[
  e_i^{c,(\ell)}=e_i^{r,(\ell+1)}\qquad(\ell\in[k-1]),
\]
contracting the resulting chain, and taking the row edges
$(e_i^{r,(1)})_{i\in[N]}$ and column edges $(e_i^{c,(k)})_{i\in[N]}$ of the chain
as the row and column edges of the result. Its symmetric matricization is
$\matA^{k}$, where $\matA=\operatorname{SymMat}(\tenA)$.

Closing the $k$-th power of $\tenA$ by further identifying
$e_i^{c,(k)}=e_i^{r,(1)}$ for every $i\in[N]$ and contracting the resulting
network with no dangling edges yields the scalar $\trace(\matA^{k})$.
\end{definition}

\begin{definition}[$(\varepsilon,\delta)$-accurate embedding]\label[definition]{def:eps-delta}
\leanref{TensorTreeSketching.IsAccurateEmbedding}\leanok
A random matrix $\matS\in\mathbb{F}^{m\times n}$ is an
\emph{$(\varepsilon,\delta)$-accurate embedding} if for every
$\boldsymbol{\vx}\in\mathbb{F}^{n}$,
\[
  \Pr_{\matS}\Bigl[\,
    \bigl|\,\TnormS{\matS\boldsymbol{\vx}}
      -\TnormS{\boldsymbol{\vx}}\,\bigr|
    >\varepsilon\,\TnormS{\boldsymbol{\vx}}
  \,\Bigr]<\delta .
\]
\end{definition}
\begin{definition}[Schatten norm]\label[definition]{def:schatten}
For a symmetric PSD matrix $\matA\in\R^{n\times n}$ with eigenvalues
$\lambda_1,\dots,\lambda_n\ge0$ and $p\ge1$, the \emph{Schatten-$p$ norm} of
$\matA$ is
\[
  \Xnorm{\matA}{p}:=\Big(\sum_{i=1}^{n}\lambda_i^{p}\Big)^{1/p}
  =\trace\bigl(\matA^{p}\bigr)^{1/p}.
\]
\end{definition}


%--------------------------------------------------------------
% Problem Statement
%--------------------------------------------------------------

\section{Problem Statement}
\label{sec:problem}
\begin{figure}[H]
\centering
\includegraphics[width=0.3\textwidth]{Illustrations/sketched_mpo_e2.png}
\caption{An example for estimator $E_2$ created by algorithm \ref{alg:mpop}} with $k=2,N=4$.
The contraction path of an MPO that contracts from left to right leaves us with this tree structure sketching tensor network.
\label{fig:mpo-e2}
\end{figure}

Let $\bigl(G_{\tenA},(e_i^{r})_{i\in[N]},(e_i^{c})_{i\in[N]},T_0\bigr)$ be a
symmetric tensor network (\Cref{def:symtn}) whose symmetric matricization
$\matA=\operatorname{SymMat}(\tenA)\in\R^{d^{N}\times d^{N}}$ is positive semidefinite matrix, and fix a power $k\ge 1$.
We draw $k$ independent samples from a sketching tensor network $G_{S}=(V_{S},E_{S})$ (\Cref{def:sketching-tn}),
denoting the $i$-th sample by $G_{S}^{(i)}$ and its matricization by
\begin{equation}
    \matS^{(i)}:=\mathrm{mat}\left(G_{S}^{(i)}\right)\in\R^{m\times d^{N}}\label{S_mat_def}
\end{equation}
Construct a contraction path on $G=(V_{\tenA}\cup V_{S},E_{\tenA}\cup E_{S})$
that evaluates
\begin{equation}\label{eq:estimator}
  E_{k}
    :=\trace\!\left(\matS^{(1)}\matA\,\matS^{(2)\top}
      \matS^{(2)}\matA\cdots
      \matS^{(k)}\matA\,\matS^{(1)\top}\right)
    \;\approx\;\trace(\matA^{k}) .
\end{equation}
We create a sketching tensor network with the following two properties.

\begin{enumerate}
  \item \textbf{$(\varepsilon,\delta)$ accuracy.} For a given $\varepsilon\in(0,1]$ and
  $\delta\in(0,1)$,
  \begin{equation}\label{eq:goal}
    \Pr\Bigl[\,\bigl|E_{k}-\trace(\matA^{k})\bigr|
      \le\varepsilon\,\trace(\matA^{k})\,\Bigr]\;\ge\;1-\delta .
  \end{equation}
  \item \textbf{Efficient evaluation.} The total amount of FLOP's and memory needed for running an algorithm to compute $E_k$  is polynomial in all system parameters.
\end{enumerate}

The rest of the paper is organized as the following: in \Cref{sec:theory}, we bound the sketch dimension $m$ needed for Eq. \eqref{eq:goal} and classify what families of tensor networks sketches satisfy  Eq. \ref{eq:goal}. In \Cref{sec:algorithms} we give the algorithm for a general symmetric tensor network and demonstrate an explicit use case for the power of a matrix product operator (MPOP) using two different algorithms that satisfy Eq. \eqref{eq:goal} and are efficiently computed.  In \Cref{sec:experiments} we report the results of our numerical experiments. 

%--------------------------------------------------------------
% Theory
%--------------------------------------------------------------
\section{Theory (temp title)}
\label{sec:theory}

\subsection{Symmetric Tensor Networks as Matrices}
\label{sec:matrixview}
We look at sketching tensor networks that have a graph network, i.e. not hyper-graphs, in such cases we can look at the estimator (Eq. \eqref{eq:estimator}) through matricizations {\color{red} [cite solomonik]}.
In that case, the analysis can be carried out entirely at the matrix level.
Given a sketching tensor network $  \bigl(G_S,\;\mu,\;m,\;(e_i^{s})_{i\in[N]}\bigr)$ where $G_S$  is a graph we can look at $$S=SymMat(G_s)$$
Throughout this section we write $d_i:=\dim(e_i^{r})=\dim(e_i^{c})$ for
$i\in[N]$ and
\[
  d_{\Pi}\;:=\;\prod_{i=1}^{N} d_i ,
\]
so that $\matA=\operatorname{SymMat}(\tenA)\in\R^{d_{\Pi}\times d_{\Pi}}$ and $\matS^{(r)}\in\R^{m\times d_{\Pi}}$. 
Note that for the case where all $d_i$'s are equal to $d$, we have $d_{\prod} =d^N$.
With this notation, This leads us to the following theorem.

{\color{blue} Leia: if this is a thrm originally from solomonik, we need to write exactly which one (number) for the cite}
\begin{theorem}[{\cite[]{masolomonik2022}}]\label{prop:ma}
\leanref{TensorTreeSketching.masolomonik_embedding}\leanaxiom
Let  $  \bigl(G_S,\;\mu,\;m,\;(e_i^{s})_{i\in[N]}\bigr)$ be a graph-structured sketching tensor network with the distribution $\mu$  either Gaussian or Rademacher with mean $0$ and standard deviation $1/\sqrt{m}$ and $N_E=|V_S|$ nodes, in which every inner edge and every dangling edge not attached to the data has size $m$. If
\[
  m=\Omega\!\bigl(N_E\log(1/\delta)/\varepsilon^{2}\bigr),
\]
then $\matS=\mathrm{mat}(G_S)$ is an $(\varepsilon,\delta)$-accurate embedding (\Cref{def:eps-delta}).
\end{theorem}

\begin{theorem}
\label{thm:k1}
\leanref{TensorTreeSketching.trace_estimation_k1}\leanok
\leanuses{trace_sketch_bound_of_eigen, IsAccurateEmbedding}
Let $\matA\in\R^{d_{\Pi}\times d_{\Pi}}$ be a symmetric PSD matrix and let $\matS\in\R^{m\times d_{\Pi}}$ be an $(\varepsilon,\delta')$-accurate embedding (\Cref{def:eps-delta}) with $\delta'=\delta/d_{\Pi}$.
Then, with probability at least $1-\delta$,
\begin{equation}
      (1-\varepsilon)\trace(\matA)\le\trace(\matS\matA\matS^{\top})
    \le(1+\varepsilon)\trace(\matA).\label{trace_eps_err}
\end{equation}
\end{theorem}

\begin{proof}
Given an $\varepsilon,\delta\in (0,1),$ and a PSD matrix $\matA$. We have that $\matA$ has the eigendecomposition
$\matA=\matU\matLambda\matU^{\top}$, with orthonormal eigenvectors $\vu_{1},\dots,\vu_{d_{\Pi}}$ and $\matLambda=\diag(\lambda_{1},\dots,\lambda_{d_{\Pi}})$, $\lambda_{i}\ge0$.
Writing $\matLambda=\matLambda^{1/2}\matLambda^{1/2}$ and $\matB:=\matU\matLambda^{1/2}$ gives $\matA=\matB\matB^{\top}$, whose columns are $\sqrt{\lambda_{i}}\,\vu_{i}$.
Now assume that we have an $\matS\in \R^{m\times d_{\Pi}}$ that for every $u_i,i=1,...,d_{\Pi}$ has the following property \begin{equation}
    (1-\varepsilon)\TnormS{u_i}\le\TnormS{Su_i}\le(1+\varepsilon)\TnormS{u_i}\label{S_eps_error}
\end{equation}
By $\trace(\matC\matC^{\top})=\FnormS{\matC}$, applied to $\matC=\matS\matB$ we have
\begin{equation}
  \trace(\matS\matA\matS^{\top})
  =\trace\!\bigl((\matS\matB)(\matS\matB)^{\top}\bigr)
  =\FnormS{\matS\matB}.
\end{equation}

Similarly, for $C=A$ we obtain 
\begin{equation}
\trace(A)=\FnormS{B}.    
\end{equation}
Summing over the columns of $\matB$ and discarding the vanishing terms, we obtain
\begin{equation}\label{eq:SB-colsum}
  \FnormS{\matS\matB}
  =\sum_{i=1}^{d_{\Pi}}\TnormS{\matS\bigl(\sqrt{\lambda_{i}}\,\vu_{i}\bigr)}
  =\sum_{i=1}^{d_{\Pi}}\lambda_{i}\TnormS{\matS\vu_{i}}\le(1+\varepsilon)\sum_{i=1}^{d_{\Pi}}\lambda_{i}\TnormS{\vu_{i}}=(1+\varepsilon)\FnormS{\matB}.
\end{equation}
Where in the last step we used Eq. \eqref{S_eps_error}. Similarly, we obtain  $$(1-\varepsilon)\FnormS{B}\le\FnormS{SB}$$
Thus, if $\matS$ has property of Eq. \eqref{S_eps_error} we get
$$  (1-\varepsilon)\trace(\matA)\le\trace(\matS\matA\matS^{\top})
    \le(1+\varepsilon)\trace(\matA)$$

If $\matS$ has the $(\varepsilon,\delta/d_{\Pi})$-accurate
property then, since the eigenvectors $\vu_1,\dots,\vu_{d_{\Pi}}$
are determined by $\matA$ alone and are therefore fixed prior to drawing
$\matS$, each of them fails to satisfy Eq. \eqref{S_eps_error} with
probability less than $\delta/d_{\Pi}$. A union bound over the $d_{\Pi}$
eigenvectors shows that with probability at least $1-\delta$ none of them
fails, hence Eq. \eqref{S_eps_error} holds simultaneously for all and Eq.
\eqref{trace_eps_err} follows.
\end{proof}

\begin{corollary}\label{cor:k1-gaussian}
\leanref{TensorTreeSketching.trace_estimation_k1_gaussian}\leanok
\leanuses{trace_estimation_k1, masolomonik_embedding}
Let $\bigl(G_S,\mu,m,(e_i^{s})_{i\in[N]}\bigr)$ be a graph-structured
sketching tensor network with $\mu$ Gaussian or Rademacher with mean $0$ and standard deviation $1/\sqrt{m}$, $N_E=|V_S|$
nodes, and every inner edge and every dangling edge not attached to the
data of size $m$. Let $\matA\in\R^{d_{\Pi}\times d_{\Pi}}$ be a symmetric
PSD matrix, and let $\varepsilon\in(0,1]$, $\delta\in(0,1)$. If
\[
  m\;=\;\Omega\!\left(\frac{N_E\log\!\bigl(d_{\Pi}/\delta\bigr)}{\varepsilon^{2}}\right)
\]

then $\matS=\mathrm{mat}(G_S)$ satisfies, with probability at least
$1-\delta$,
\[
  (1-\varepsilon)\trace(\matA)
  \;\le\;\trace\!\bigl(\matS\matA\matS^{\top}\bigr)
  \;\le\;(1+\varepsilon)\trace(\matA) .
\]
\end{corollary}

\begin{proof}
Apply \Cref{prop:ma} with failure probability $\delta'=\delta/d_{\Pi}$: the
requirement $m=\Omega\bigl(N_E\log(1/\delta')/\varepsilon^{2}\bigr)$ is exactly the
stated bound, so $\matS$ is an $(\varepsilon,\delta')$-accurate embedding. The
conclusion follows from \Cref{thm:k1}.
\end{proof}

\subsection{The case \texorpdfstring{$k\ge2$}{k>=2}}
\label{sec:tracek}
We now generalize for higher powers, i.e. $k>1$. 
\begin{theorem}\label{thm:tracek}
\leanref{TensorTreeSketching.trace_estimation_k}\leanok
\leanuses{gram_control, deviation_bound_rel}
Let $\matA\in\R^{d_{\Pi}\times d_{\Pi}}$ be a symmetric PSD matrix, $k \ge 2$, $\hat\epsilon \in (0,1]$, $\delta \in (0,1)$, denote $\rho :=\frac{\trace(\matA)}{\Xnorm{\matA}{k}}$ (\Cref{def:schatten}) and let each $\matS^{(r)}, r\in[k]$ have the $(\epsilon,\delta')$-embedding property with
\begin{equation}\label{eq:eps-choice}
\delta' \;=\; \frac{\delta}{k\, d_{\Pi}^2},
\qquad
\epsilon \;=\; \min\Big\{ \frac{\hat\epsilon}{2k},\ \frac{1}{k\rho}\sqrt{\hat\epsilon/2} \Big\}.
\end{equation}
Then with the definition for $E_k$ in Eq. \eqref{eq:estimator} we have
\begin{equation}\label{eq:thm-conclusion}
\mathbb{P}\Big( \big| E_k - \trace(\matA^k) \big| \;\le\; \hat\epsilon\, \trace(\matA^k) \Big) \;\ge\; 1-\delta .
\end{equation}
\end{theorem}
Next, we use this theorem for a sketching tensor network algorithm. We find the desired property for the required accuracy with the following corollary.
\begin{corollary}\label[corollary]{cor:k-gaussian}
\leanref{TensorTreeSketching.trace_estimation_k_gaussian}\leanok
\leanuses{trace_estimation_k, masolomonik_embedding}
Let $\bigl(G_S,\mu,m,(e_i^{s})_{i\in[N]}\bigr)$ be a graph-structured sketching tensor network with $\mu$ Gaussian or Rademacher with mean $0$ and standard deviation $1/\sqrt{m}$, $N_E=|V_S|$ nodes, and every inner edge and every dangling edge not attached to the data of size $m$.
Let $\matA\in\R^{d_{\Pi}\times d_{\Pi}}$ be a symmetric PSD matrix, let $k\ge2$, $\hat\varepsilon\in(0,1]$, $\delta\in(0,1)$,
and $\rho=\trace(\matA)/\Xnorm{\matA}{k}$ (\Cref{def:schatten}). 
If
\[
  m\;=\;\Omega\!\left(
    \frac{N_E\,k^{2}\max\{1,\rho^{2}\hat\varepsilon\}}{\hat\varepsilon^{2}}
    \bigl(\log d_{\Pi}+\log(k/\delta)\bigr)\right),
\]
then the estimator $E_k$  (\ref{eq:estimator}) built from $k$ independent samples $\matS^{(1)},\dots,\matS^{(k)}$ of $G_S$ satisfies, with probability at least $1-\delta$,

\[
  (1-\hat\varepsilon)\,\trace(\matA^{k})
  \;\le\; E_{k} \;\le\;
  (1+\hat\varepsilon)\,\trace(\matA^{k}) .
\]
\end{corollary}
\begin{proof}
Set $\delta'=\delta/(k\,d_{\Pi}^2)$ and $\epsilon=\min\bigl\{\hat\varepsilon/(2k),\,(k\rho)^{-1}\sqrt{\hat\varepsilon/2}\bigr\}$ as in Eq. \eqref{eq:eps-choice}.
Up to constant factors, the stated bound on $m$ is therefore the requirement $m=\Omega\bigl(N_E\log(1/\delta')/\epsilon^{2}\bigr)$ of \Cref{prop:ma}, so each $\matS^{(r)}$ is an $(\epsilon,\delta')$-accurate embedding. The claim now follows from \Cref{thm:tracek}.
\end{proof}


\begin{lemma}[Two-sided control of the sketched Gram matrices]\label[lemma]{lem:gram}
\leanref{TensorTreeSketching.gram_control}\leanok
Given a PSD matrix $\matA\in\R^{d_{\Pi}\times d_{\Pi}}$ and then let $\matA=\matU\matLambda\matU^{\top}$ be its eigen-decomposition, with orthonormal eigenvectors $u_{1},\dots,u_{d_{\Pi}}$ and $\matLambda=\diag(\lambda_{1},\dots,\lambda_{d_{\Pi}})$, and set
\[
  (\Gamma_r)_{ij} \;:=\; \inner{\matS^{(r)} u_i}{\matS^{(r)} u_j} - \delta_{ij},
  \qquad r\in[k],\ i,j\in[d_{\Pi}].
\]
If each $\matS^{(r)}$ has the $(\epsilon,\delta')$-embedding property, then with probability at least $1 - k\, d_{\Pi}^2 \delta'$, simultaneously,
\begin{equation}\label{eq:gram-event}
\big| (\Gamma_r)_{ij} \big| \;\le\; \epsilon
\qquad \text{for all } r \in [k],\ i,j \in [d_{\Pi}].
\end{equation}
\end{lemma}

\begin{proof}
Fix $r$. For $i=j$: $(\Gamma_r)_{ii} = \TnormS{\matS^{(r)} u_i} - 1$ and $\TnormS{u_i} = 1$, so $|(\Gamma_r)_{ii}| \le \epsilon$ with probability $\ge 1-\delta'$. For $i \neq j$: the orthonormality gives $\TnormS{u_i \pm u_j} = 2$, and by the polarization identity together with $\inner{u_i}{u_j} = 0$, we obtain
\begin{align}
(\Gamma_r)_{ij} \;=\; \inner{S_r u_i}{S_r u_j}
&= \tfrac14\Big( \big(\TnormS{S_r(u_i+u_j)} - 2\big) \;-\; \big(\TnormS{S_r(u_i-u_j)} - 2\big) \Big), \\
\big|(\Gamma_r)_{ij}\big| &\le \tfrac14\big( 2\epsilon + 2\epsilon \big) \;=\; \epsilon
\end{align}
on the intersection of the two embedding events for the fixed vectors $u_i + u_j$ and $u_i - u_j$. Per sketch this is a union over $d_{\Pi} + 2\binom{d_{\Pi}}{2} = d_{\Pi}^2$ fixed vectors, hence over $k\, d_{\Pi}^2$ events across the $k$ sketches, each of probability at most $\delta'$.
\end{proof}

\begin{lemma}[Deterministic deviation bound]\label[lemma]{lem:master}
\leanref{TensorTreeSketching.deviation_bound_rel}\leanok
\leanuses{deviation_bound_abs, trPow_le_specNorm_mul, one_add_pow_sub_le}
Let $\epsilon \ge 0$ and suppose $|(\Gamma_r)_{ij}| \le \epsilon$ for all $r \in [k]$, $i,j \in [d_{\Pi}]$. For $\emptyset \neq T = \{r_1 < \dots < r_j\} \subseteq [k]$ set $q_i := r_{i+1} - r_i$ for $1 \le i < j$ and $q_j := k - r_j + r_1$, so that $q_i \ge 1$ and $\sum_{i=1}^{j} q_i = k$. Then
\begin{equation}\label{eq:master}
\big| E_k - \operatorname{Tr}(\matA^k) \big|
\;\le\;
\sum_{\emptyset \neq T \subseteq [k]} \epsilon^{|T|} \prod_{i=1}^{|T|} \operatorname{Tr}\!\big(\matA^{q_i}\big)
\;\le\;
\big( \Tnorm{\matA} + \epsilon \operatorname{Tr}(\matA) \big)^k - \Tnorm{\matA}^k ,
\end{equation}
dividing by $\operatorname{Tr}(\matA^k)$,\begin{equation}\label{eq:master-rel}
\frac{\big| E_k - \operatorname{Tr}(\matA^k) \big|}{\operatorname{Tr}(\matA^k)}
\;\le\;
k\epsilon \;+\; \sum_{j=2}^{k} \binom{k}{j} (\epsilon\rho)^j
\;\le\;
k\epsilon + (k\epsilon\rho)^2
\qquad \text{whenever } k\epsilon\rho \le 1 .
\end{equation}
\end{lemma}

\begin{proof}
Write $M_r := \matU^{\top}\matS^{(r)\top}\matS^{(r)}\matU = I + \Gamma_r$ . By cyclicity of the trace, \eqref{eq:estimator} equals $\trace\big(\prod_{r=1}^{k} \matLambda M_r\big)$.
expanding multilinearly:
\begin{equation}\label{eq:expansion}
E_k \;=\; \operatorname{Tr}\!\Big( \prod_{r=1}^{k} (I + \Gamma_r) \Lambda \Big)
\;=\; \sum_{T \subseteq [k]} \tenT_T,
\qquad
\tenT_\emptyset = \operatorname{Tr}(\Lambda^k) = \operatorname{Tr}(\matA^k).
\end{equation}
For $T = \{r_1 < \dots < r_j\}$, collapsing the identity factors merges consecutive powers of $\Lambda$, and cyclicity of the trace places exactly $q_i$ factors of $\Lambda$ between $\Gamma_{r_i}$ and $\Gamma_{r_{i+1}}$:
\begin{equation}\label{eq:term-index}
\tenT_T
\;=\; \operatorname{Tr}\!\big( \Gamma_{r_1} \Lambda^{q_1} \Gamma_{r_2} \Lambda^{q_2} \cdots \Gamma_{r_j} \Lambda^{q_j} \big)
\;=\; \sum_{a_1,\dots,a_j = 1}^{d_{\Pi}}\ \prod_{i=1}^{j} (\Gamma_{r_i})_{a_i a_{i+1}}\, \lambda_{a_{i+1}}^{\,q_i},
\qquad a_{j+1} := a_1 .
\end{equation}
Given some subset $\emptyset \neq T \subseteq [k]$ with $|T| = j$, by the triangle inequality, $|(\Gamma_{r_i})_{a_i a_{i+1}}| \le \epsilon$, $\lambda_a \ge 0$, and the fact that each summation index $a_{i+1}$ occurs in exactly one eigenvalue factor,

 
\begin{equation}\label{eq:perterm}
|\tenT_T| \;\le\; \epsilon^{j} \sum_{a_1,\dots,a_j=1}^{d_{\Pi}} \prod_{i=1}^{j} \lambda_{a_{i+1}}^{\,q_i}
\;=\; \epsilon^{j} \prod_{i=1}^{j} \Big( \sum_{a=1}^{d_{\Pi}} \lambda_a^{\,q_i} \Big)
\;=\; \epsilon^{j} \prod_{i=1}^{j} \operatorname{Tr}\!\big(\matA^{q_i} \big).
\end{equation}
Summing over $T \neq \emptyset$ yields the first inequality in Eq. \eqref{eq:master}. Since $\lambda_a^{\,q} \le \Tnorm{\matA}^{\,q-1} \lambda_a$,
\begin{equation}\label{eq:tr-power}
\operatorname{Tr}(\matA^{q}) \;\le\; \Tnorm{\matA}^{\,q-1} \operatorname{Tr}(\matA),
\qquad \text{hence} \qquad
\prod_{i=1}^{j} \operatorname{Tr}(\matA^{q_i}) \;\le\; \Tnorm{\matA}^{\,k-j}\, \big(\operatorname{Tr} \matA\big)^{j},
\end{equation}
and there are $\binom{k}{j}$ subsets of size $j$, so the middle term of Eq. \eqref{eq:master} is at most
$\sum_{j=1}^{k} \binom{k}{j} \epsilon^j (\operatorname{Tr}\matA)^j \Tnorm{\matA}^{\,k-j} = (\Tnorm{\matA} + \epsilon\operatorname{Tr}\matA)^k - \Tnorm{\matA}^k$.

For Eq. \eqref{eq:master-rel}: if $|T| = 1$ then $q_1 = k$ and Eq. \eqref{eq:perterm} reads $|\tenT_T| \le \epsilon \operatorname{Tr}(\matA^k)$; there are $k$ such terms. 
For $j \ge 2$: applying $\Tnorm{\matA}\le\Xnorm{\matA}{k}$
(\Cref{def:schatten}) to Eq. \eqref{eq:tr-power} gives
$\operatorname{Tr}(\matA^{q_i}) \le \Xnorm{\matA}{k}^{\,q_i-1}\operatorname{Tr}(\matA)$;
multiplying over $i$, with $\sum_{i=1}^{j}(q_i-1)=k-j$ and
$\Xnorm{\matA}{k}^{k}=\operatorname{Tr}(\matA^{k})$,
\begin{equation}
\epsilon^{j} \prod_{i=1}^{j} \operatorname{Tr}(\matA^{q_i})
\;\le\; \epsilon^{j}\, \big(\operatorname{Tr}\matA\big)^{j}\, \Xnorm{\matA}{k}^{\,k-j}
\;=\; (\epsilon\rho)^{j}\, \operatorname{Tr}(\matA^k)
\end{equation}

\
Finally, for $x := \epsilon\rho$ with $kx \le 1$,
\begin{equation}
\sum_{j=2}^{k} \binom{k}{j} x^j \;=\; (1+x)^k - 1 - kx \;\le\; e^{kx} - 1 - kx \;\le\; (kx)^2 ,
\end{equation}
using $e^t \le 1 + t + t^2$ for $t \in [0,1]$.
\end{proof}

\begin{proof}[Proof of Theorem~\ref{thm:tracek}]
By Lemma~\ref{lem:gram} and the choice $\delta' = \delta/(k d_{\Pi}^2)$, the event of Eq. \eqref{eq:gram-event} holds with probability at least $1-\delta$. On this event apply Lemma~\ref{lem:master}: the choice from Eq. \eqref{eq:eps-choice} gives $k\epsilon \le \hat\epsilon/2$ and $(k\epsilon\rho)^2 \le \hat\epsilon/2 \le 1$ (in particular $k\epsilon\rho \le 1$), hence by Eq. \eqref{eq:master-rel}
\begin{equation}
\big| E_k - \operatorname{Tr}(\matA^k) \big| \;\le\; \Big( \frac{\hat\epsilon}{2} + \frac{\hat\epsilon}{2} \Big) \operatorname{Tr}(\matA^k) \;=\; \hat\epsilon\, \operatorname{Tr}(\matA^k).
\end{equation}
\end{proof}

%--------------------------------------------------------------
% Algorithms
%--------------------------------------------------------------
\section{Algorithms}
\label{sec:algorithms}
We now turn the guarantees of \Cref{sec:theory} into concrete procedures. In
every case the input is a symmetric tensor network
$\bigl(G_{\tenA},(e_i^{r}),(e_i^{c}),T_0\bigr)$ and a power $k$, and the
output is the scalar estimate $E_{k}\approx\trace(\matA^{k})$ of Eq.
\eqref{eq:estimator} where the $(\varepsilon,\delta)$ accuracy is inherited from
\Cref{thm:k1,thm:tracek}.
We start with introducing a general framework for a general symmetric tensor network in subsection \ref{sec:alg-general}, then we give 2 possible sketches for the explicit use case for the MPOP case in subsections \ref{sec:alg-mpo} and \ref{sec:alg-mpo-2}.

\subsection{Sketching Tensor Network Family}

\begin{figure}[H]
    \centering
    \includegraphics[width=0.5\linewidth]{Illustrations/guy_files/three_sketches.png}
    \caption{Sketch Tensor Networks Examples}
    \label{fig:sketch_examples}
\end{figure}

From theorems in \ref{thm:k1} and \ref{thm:tracek} we get that the requirements are fulfilled with many types tensor networks. For a sketching tensor network $G_S$ that is a graph, we only require that all edges not connected to the data tensor network $G_A$  be of dimension $m$. This requirement can be fulfilled by several different graph tensor networks (See Fig. \ref{fig:sketch_examples}). 
For our calculations to be efficient, we choose a graph with a single sketch dimension, as explained in subsection \ref{sec:alg-general}, which leaves us with fewer possibilities.


\subsection{A general tensor-network framework}
\label{sec:alg-general}

For a general symmetric tensor network ($G_\tenA$, $(e_i^{r})_{i\in[N]}$, $(e_i^{c})_{i\in[N]}$, $T_0$) (\Cref{def:symtn}), and a given $k\in\N$. We use \Cref{alg:trace} to efficiently calculate our estimator from Eq. \eqref{eq:estimator} at the tensor network level. 
We start by sketching each tensor network $G_\tenA$ to a matrix shaped $m\times m$, we do this by sketching the rows and columns dimensions with a sketch tensor networks $G_S$  that has a  single output sketch dimension. This gives us a matrix shaped $m\times m$  then we work at the matrix level by multiplying each matrix and calculating the trace of the final matrix.

For a general symmetric tensor network $\bigl(G_{\tenA},(e_i^{r})_{i\in[N]},(e_i^{c})_{i\in[N]},T_0\bigr)$ and $k\in\N$, \Cref{alg:trace} evaluates the estimator (Eq. \eqref{eq:estimator}) at the tensor network level. The sketch is applied during the data contraction path rather than formed explicitly. Whenever a contraction of $T_0$ would leave two dangling row (respectively column) edges on its output, a sketching random tensor merges them into a single sketch leg of size $m$. Every such merge realizes a unique node of the common dimension tree $T_D$, and the tensor placed there is keyed by that dimension tree node together with a sample index. It is drawn at its first occurrence and reused at its second. This realizes the shared randomness of Eq. \eqref{eq:estimator}. Each replica thus contracts to a matrix $\matM_i=\matS^{(i)}\matA\matS^{(i+1)\top}\in\C^{m\times m}$, and $E_k=\trace(\matM_1\cdots\matM_k)$.

\paragraph{Node sketch tensors.}
For each internal node $t\in T_D$  for $T_D$ being the dimension tree (\Cref{def:dimtree}), and each sample index $s\in[k]$ we associate an
order-$3$ tensor $\tenS_t^{(s)}$ with i.i.d.\ Gaussian (or Rademacher) with mean $0$ and standard deviation $1/\sqrt{m}$ entries. Two legs match the two dangling legs merged at $t$, and the third, of size $m$, replaces them. Each $\tenS_t^{(s)}$ is drawn at its first use and reused afterwards, with the cyclic identification
$\tenS_t^{(k+1)}:=\tenS_t^{(1)}$.
\begin{algorithm}[H]
\caption{Sketched estimation of $\trace(\matA^{k})$ (general symmetric network)}
\label{alg:trace}
\begin{algorithmic}[1]
\STATE \textbf{Input:} symmetric tensor network $\bigl(G_{\tenA},(e_i^{r}),(e_i^{c}),T_0\bigr)$ with common dimension tree $T_D$; power $k$; sketch size $m$.
\STATE \textbf{Output:} estimate $E_k$ of $\trace(\matA^{k})$.
\STATE \emph{Reshape} co-located dangling legs, so every vertex carries at most one dangling row leg and one dangling column leg.
\FOR{$r=1,\dots,k$}
  \FOR{each contraction $(U_i,V_i)$ of $T_0$ on a fresh replica of $G_{\tenA}$}
    \IF{both inputs carry a dangling row leg, realizing node $t\in T_D$}
      \STATE connect a sketching tensor $\tenS_t^{(r)}$ to the two row legs
    \ENDIF
    \IF{both inputs carry a dangling column leg, realizing node $t'\in T_D$}
      \STATE connect a sketching tensor $\tenS_{t'}^{(r+1)}$ to the two column legs
    \ENDIF
    \STATE contract the resulting sub tensor network of $2$--$4$ nodes, choosing the cheapest contraction order
  \ENDFOR
  \STATE the replica contracts to $\matM_r\in\R^{m\times m}$
\ENDFOR
\STATE $\matM\gets\matM_1\matM_2\cdots\matM_k$
\STATE $E_k\gets\trace(\matM)$
\RETURN $E_k$
\end{algorithmic}
\end{algorithm}

\begin{lemma}[Correctness of \Cref{alg:trace}]\label[lemma]{lem:correct}
% Matrix-level shadows of this correctness argument (former lem:correct + cor:chain):
% each replica contracts to M_r = S_r A S_{r+1}^T, and their product's trace is E_k.
\leanref{TensorTreeSketching.contract2_eq_mul}\leanok\ \leanref{TensorTreeSketching.chained_copies_eq_estimator}\leanok
\leanuses{contract2_assoc}
The output of \Cref{alg:trace} is exactly the estimator $E_{k}$ of Eq.
\eqref{eq:estimator} and the contraction tree executed by \Cref{alg:trace} is constrained on
$T_0$ (\Cref{def:constrained}). In particular, for $m$ as in \Cref{cor:k1-gaussian} ($k=1$) or \Cref{cor:k-gaussian} ($k\ge2$), it satisfies
the $(\varepsilon,\delta)$ guarantee \eqref{eq:goal}.
\end{lemma}
\begin{proof}

We count what the algorithm does on a single replica. The
contraction path $T_0$ consists of $|V_{\tenA}|-1$ contractions. The
dimension tree $T_D$ (\Cref{def:dimtree}) has $N$ leaves - the
singletons $\{e_i\}$, $i\in[N]$ - and, being a full binary tree,
exactly $N-1$ internal nodes; we denote them $t_1,\dots,t_{N-1}$, each
a subset of $[N]$ formed by merging its two children. Since the row and
the column dimension trees of $T_0$ are both equal to $T_D$
(\Cref{def:symtn}), each node $t_i$ is realized by exactly two of the
$|V_{\tenA}|-1$ contractions: one merging the operands that carry the
row edges of its children, and one merging those that carry their
column edges. These $2(N-1)$ contractions are precisely the ones at
which \Cref{alg:trace} places sketch tensors; all remaining
contractions involve data tensors only.

The sketch tensors are $\tenS_i^{(r)}$ for $i\in[N-1]$ and $r\in[k]$
(cyclically, $\tenS_i^{(k+1)}:=\tenS_i^{(1)}$): the tensor
$\tenS_i^{(r)}$ is the order-3 tensor of sample $r$ associated with the
node $t_i$, with two input legs keyed by the two children of $t_i$ and
one output leg of size $m$. On replica $r$, the algorithm places
$\tenS_i^{(r)}$ at the row-side contraction realizing $t_i$ and
$\tenS_i^{(r+1)}$ at the column-side one.

Since the input legs are keyed by the children of $t_i$, connecting a
sketch tensor leaves no freedom: each leg meets the data legs (if the
child is a leaf) or the output leg (if the child is a node $t_j$, whose
tensor was placed earlier along $T_0$) of \emph{its} child - on the
row side through the child's row edges, and on the column side, for the
same leg of the tensor with the same key, through the child's column
edges. The correspondence between the two sides is well defined because
the two dimension trees are equal \emph{as trees labeled by edge-index
sets}. Consequently the row-side placements of replica $r$ assemble the
sample $G_S^{(r)}$ of the sketching tensor network built on $T_D$,
attached to $(e_i^{r})_{i\in[N]}$, and the column-side placements
assemble the sample $G_S^{(r+1)}$ - the same tensors, connected the
same way - attached to $(e_i^{c})_{i\in[N]}$
(\Cref{def:sketching-tn}). The two attachments realize
$\matS^{(r)}=\mathrm{mat}\bigl(G_S^{(r)}\bigr)$ acting on the rows and
$\matS^{(r+1)\top}$ acting on the columns, so replica $r$ contracts to
$\matM_r=\matS^{(r)}\matA\,\matS^{(r+1)\top}$, and the output
$E_k=\trace(\matM_1\cdots\matM_k)$ is exactly the estimator
\eqref{eq:estimator}.
Finally, each local contraction of nodes contains exactly one step
merging the operand containing $U_i$ with the operand containing $W_i$, and
at that step $\hat U\cap V_{\tenA}=U_i$ and $\hat W\cap V_{\tenA}=W_i$ getting
the executed contraction tree is constrained on $T_0$.
\end{proof}
\subsubsection{The matrix product operator case}
\label{sec:alg-mpo}

\paragraph{Physical origin: the reduced density matrix.}
We wish to solve a specific problem in quantum many-body physics and quantum information \cite{vidal2003,schollwock2011,oseledets2011tt}  {\color{red} + cite exactly like Raphael which sites Noa also + cite my founded Entanglement R´enyi entropies in holographic theories.}
Consider a quantum state $\vpsi\in\C^{d^{N_A}}\kronprod\C^{d^{N_{\overline{A}}}}$ describing
two interacting subsystems $A$  and $\overline{A}$ with $N_A$  and $N_{\overline{A}}$  sites respectively where we assume for simplicity that each sites local dimension is $d$. 
The state of the subsystem $A$ alone is captured by the \emph{reduced density matrix} (RDM) defined
\[
  \rho_A=\trace_{\overline{A}}(\vpsi\vpsi^{*})\in\C^{d^{N_A}\times d^{N_A}},
\]
We wish to calculate the R\'enyi moments as measures of entanglement defined as:
\begin{equation}
    R_k:=Tr(\rho_A^k)
\end{equation}
We will focus this algorithm when our state is a one-dimensional spin system with $N$  sites where each site's physical dimension is $d$. Thus, our vector is $\vpsi\in\C^{d^{N_A}}$. Avoiding storing the entire exponentially large vector in memory, we  write the system in tensor network form.
\\
A \emph{matrix product state} (MPS) represents a vector
$\vpsi\in\C^{d^{N}}$ on $N$ sites by $N-2$ order-$3$ cores
$\tenB_{i}\in\C^{\chi\times d\times\chi}$, $i=2,\dots,N-1$, and boundary
cores $\tenB_{1},\tenB_{N}\in\C^{d\times\chi}$, each core carrying one
dangling  physical index of dimension $d$; the non-dangling
\emph{bond} indices, of dimension $\chi$, are summed (see \Cref{fig:mps-mpo-mpop}(a)):
\begin{equation}\label{eq:mps-def}
  \vpsi_{(i_1,\dots,i_N)}
  =\sum_{\alpha_1,\dots,\alpha_{N-1}=1}^{\chi}
   \prod_{n=1}^{N}\tenB_{n}[\alpha_{n-1},i_n,\alpha_n].
\end{equation}.

\begin{figure}[H]
\centering
\includegraphics[width=\textwidth]{Illustrations/mps_mpo_mpop.png}
\caption{(a) an MPS with physical dimension $d$  and bond dimension $\chi$ with its complex conjugate (b) An MPO created by taking each MPS site and using a Kronecker operation site by site ; (c) the MPO power:
$k$ stacked copies. Dangling legs have dimension $d$, bonds dimension $D$. When merging column by column, we are left with the bond dimension of $D^k$.}
\label{fig:mps-mpo-mpop}
\end{figure}
A \emph{matrix product operator} (MPO) 
represents a matrix $\matA\in\C^{d^{N}\times d^{N}}$ on $N$ sites by $N-2$
order-$4$ cores $\ten{A}_i\in\C^{D\times d\times d\times D}, i=2,...,N-1$ per site and $A_1,A_N\in\C^{d\times d\times D}$, each
carrying 2 dangling dimensions with dimension $d$, called the physical index, one for the rows and one for the columns dimensions. The remaining dimensions, non-dangling edges, for each tensor are called the bond indices, and for simplicity we assume that they are dimension $D$ (see \Cref{fig:mps-mpo-mpop}(b)).
The entries of $\matA$ are
\begin{equation}\label{eq:mpo-def}
  \matA_{(i_1,\dots,i_N),(j_1,\dots,j_N)}
  \;=\;
  \sum_{\alpha_1,\dots,\alpha_{N-1}=1}^{D}
  \;\prod_{n=1}^{N}\ten{A}_n[\alpha_{n-1},i_n,j_n,\alpha_n],
\end{equation}
\\
A \emph{matrix product operator power} (MPOP) represents a matrix
$\matA^{k}\in\C^{d^{N}\times d^{N}}$ where $\matA^{k}$ is itself an MPO. The
multiplication acts site by site, at each site $n$ the column physical
index of one copy is contracted with the row physical index of the next,
while the $k$ bond indices stack into one merged bond
\cite{schollwock2011,hubig2017mpo} (see \Cref{fig:mps-mpo-mpop}(c)).

In our case, we obtain an MPO with the nodes from the MPS: represent the projector
$\vpsi\vpsi^{*}$ by stacking two copies of the MPS from Eq.\eqref{eq:mps-def} - the copy of $\vpsi$ with its physical legs facing upwards, and the
complex-conjugated copy with its physical legs facing downwards. Reading the stacked pair site by site, the two cores at each site form a single order-$4$ core with one
upward physical leg (a row index), one downward physical leg (a column
index), and the two bond pairs merged into one bond of dimension
$\chi^{2}$. Mathematically, we use a Kronecker product per site, i.e. each MPO site is comprised of $\tenA_i=\tenB_i\kronprod \tenB_i^*$.
Comparing with Eq. \eqref{eq:mpo-def}, the network
$\vpsi\vpsi^{*}$ is exactly an MPO with bond dimension $D=\chi^{2}$. The partial trace then
becomes a network operation, $\trace_{B}$ contracts each upward leg with
the downward leg beneath it on the sites of $\overline{A}$, and what remains is the
RDM $\rho_{A}$, an MPO on the $N_{A}$ sites of $A$, still of bond
dimension $\chi^{2}$. 



\subsubsection{The MPO case}
We now instantiate \Cref{alg:trace} for a tensor network given as an MPO
(\Cref{alg:mpop} below). The algorithm takes as input an MPO with $N$ sites, physical dimension $d$, and bond
dimension $D$, and computes the estimator $E_k$ of Eq. \eqref{eq:estimator}.
We adopt the trivial contraction path, the network is contracted from left
to right, i.e., $N-1$ contractions over the $N$ site columns.
\Cref{alg:trace} attaches to each such contraction two order-$3$ sketch
cores, one above and one below the MPO core, so each step contracts
four tensors.
It remains to fix the order of this four-tensor contraction, which we choose based on the parameters.
We choose our contraction order for these $4$  tensors assuming that $d$  is small, our $m$  can grow with the network parameters all the while $d=O(1)$.  With this in mind we choose the optimal contraction path for the $4$  tensors. 


\begin{algorithm}[H]
\caption{Sketched estimation of $\trace(\matA^{k})$ (MPO)}
\label{alg:mpop}
\begin{algorithmic}[1]
\STATE \textbf{Input:} MPO cores $\tenA_i$, $i\in[N]$; power $k$; sketch size $m$.
\STATE \textbf{Output:} estimate $E_k$ of $\trace(\matA^{k})$.
\STATE \emph{Sampling convention:} each sketch core $\tenS_n^{(s)}$ is drawn i.i.d.\ Gaussian or Rademacher with with mean $0$ and standard deviation $1/\sqrt{m}$ at its first occurrence and reused at its second occurrence. $\matS^{(k+1)}:=\matS^{(1)}$.\FOR{$r=1,\dots,k$}
  \STATE connect sketch cores $\tenS_1^{(r)}$ and $\tenS_1^{(r+1)}$ to the row and column legs of $\tenA_1$
 \STATE $\ten{W}\gets\bigl((\tenS_1^{(r)}\cdot\tenA_1)\cdot\tenA_2\bigr)\cdot\tenS_1^{(r+1)}$
  \FOR{$n=2,\dots,N-2$}
    \STATE connect sketch cores $\tenS_n^{(r)}$ and $\tenS_n^{(r+1)}$ to the row and column legs of $\tenA_{n+1}$
    \STATE $\ten{W}\gets\bigl((\ten{W}\cdot\tenS_n^{(r)})\cdot\tenA_{n+1}\bigr)\cdot\tenS_n^{(r+1)}$
  \ENDFOR
  \STATE $\matM_r\gets\bigl((\ten{W}\cdot\tenS_{N-1}^{(r)})\cdot\tenA_N\bigr)\cdot\tenS_{N-1}^{(r+1)}$
\ENDFOR
\STATE $\matM\gets\matM_1\matM_2\cdots\matM_k$
\STATE $E_k\gets\trace(\matM)$
\RETURN $E_k$
\end{algorithmic}
\end{algorithm}

\begin{lemma}[Correctness and cost of \Cref{alg:mpop}]\label{lem:mpop-correct}
Given an MPO \eqref{eq:mpo-def} with parameters $N,d,D$ whose matricization is
positive semidefinite, and $k\ge1$: the output of \Cref{alg:mpop} is exactly
the estimator $E_k$ of Eq. \eqref{eq:estimator} and the executed contraction tree is constrained on the
left-to-right path $T_0$ (\Cref{def:constrained}). For $m$ as in
\Cref{cor:k1-gaussian} ($k=1$) or \Cref{cor:k-gaussian} ($k\ge2$) with
$d_{\Pi}=d^{N}$, it satisfies the $(\varepsilon,\delta)$ guarantee
from Eq. \eqref{eq:goal} . \\
The total cost is $O\!\bigl(kN\,dD\,m^{3}+kN\,d^{2}D^{2}m^{2}+k\,m^{3}\bigr)$, which is polynomial in all network parameters $d,D,N,k$.
\end{lemma}
\begin{proof}The sketch assembled by \Cref{alg:mpop} (For example, Fig. \ref{fig:mpo-e2}) is a sample of the sketching tensor network (\Cref{def:sketching-tn}).
What remains is choosing the output edge of size $m$ \Cref{prop:ma}.
Hence, with $d_{\Pi}=d^{N}$, \Cref{cor:k1-gaussian} gives for $k=1$ the $(\varepsilon,\delta)$ guarantee Eq. \eqref{eq:goal}
whenever
\[
  m=\Omega\!\left(\frac{N\bigl(N\log d+\log(1/\delta)\bigr)}
    {\varepsilon^{2}}\right),
\]
and \Cref{cor:k-gaussian} gives for $k\ge2$ the guarantee whenever
\[
m=\Omega\!\left(\frac{N\,k^{2}\max\{1,\rho^{2}\varepsilon\}}
    {\varepsilon^{2}}\bigl(N\log d+\log(k/\delta)\bigr)\right),
\]

Finally, \Cref{alg:mpop} is the instantiation of \Cref{alg:trace} on the
trivial left-to-right contraction path, so by \Cref{lem:correct} each
replica contracts to $\matM_r=\matS^{(r)}\matA\,\matS^{(r+1)\top}$ and the output is exactly the estimator from Eq. \eqref{eq:estimator}, the
$(\varepsilon,\delta)$ accuracy therefore transfers unchanged. 
By the same lemma, the executed contraction tree is constrained on the
left-to-right path $T_0$ (\Cref{def:constrained}).
\leanref{TensorTreeSketching.estimator_eq_transfer_prod}\leanok
\leanuses{trace_transfer_prod_cyclic, chained_copies_eq_estimator}


\medskip\noindent\emph{Computational Cost.}
For each MPO replica the loop runs $N-2$ times, each iteration contracting
its 4 tensors at cost $2dD\,m^{3}+d^{2}D^{2}m^{2}$
(\Cref{rem:contraction-cost}) and the first and last contractions, with the order-$3$
boundary core, cost $O(dD\,m^{2})$. Getting that a
single sketched MPO costs $O\!\bigl(N(dD\,m^{3}+d^{2}D^{2}m^{2})\bigr)$.
Repeating $k$ times and then multiply the resulting $m\times m$ matrices at a
cost of $k\,m^{3}$, and take the final trace at cost $m$. The total computational cost described is $$O\!\bigl(kN\,dD\,m^{3}\;+\;kN\,d^{2}D^{2}m^{2}\;+\;k\,m^{3}\bigr)$$
i.e. polynomial in all system parameters, fulfilling the efficiency
requirement of \Cref{sec:problem}.
\end{proof}

\subsubsection{The MPS case}
\label{sec:alg-mpo-2}

\paragraph{MPS Algorithm}
We now take it a step further and exploit more of the structure of a tensor network, as we treat the the input as an MPS $\tenB$ on $N$ sites
(Eq. \eqref{eq:mps-def}), representing the vector $\vb\in\R^{d^{N}}$, and the
operator is the rank-one projector $\matA=\vb\vb^{\top}$. Rather than forming the MPO of $\matA$ and running
\Cref{alg:mpop}, we sketch the factors instead. 
Each replica block factorizes as
\[
  \matM_{r}
  =\matS^{(r)}\matA\,\matS^{(r+1)\top}
  =\bigl(\matS^{(r)}\vb\bigr)\bigl(\matS^{(r+1)}\vb\bigr)^{\top}
  =\vv_{r}\vv_{r+1}^{\top},
  \qquad
  \vv_{r}:=\matS^{(r)}\vb\in\C^{m},
\]
so the estimator collapses to $k$ dot products,
\[
  E_{k}
  =\trace\bigl(\vv_{1}\vv_{2}^{\top}\,\vv_{2}\vv_{3}^{\top}\cdots
    \vv_{k}\vv_{1}^{\top}\bigr)
  =\prod_{r=1}^{k}\inner{\vv_{r}}{\vv_{r}}.
\]
We note that the tensor network $\tenB$  is not a symmetric tensor network. We can leverage known algorithms for a better optimal cost, as long as the matrix $\matS$ obtained from such algorithm supports our theory.
To compute the sketched vectors $\vv_{r}$ we use the sketching algorithm
of \cite{masolomonik2022}, which attains a lower asymptotically optimal sketching cost. 
\\
Their sketch in \cite{masolomonik2022} is applied in two parts. First, a per-node Kronecker step: the physical leg of each site
is sketched by an independent $(d,m)$ tensor, leaving $N$ sketched nodes
$\hat{\tenB}_{n}$, swapping the $d$  physical dimension with $m$.
Then, along the left-to-right contraction path, each
contraction $i$ of copy $r$ receives two gadget tensors
$g_{i}^{(r,a)}\in\R^{m\times\sqrt{m/\chi}\times m}$ and
$g_{i}^{(r,b)}\in\R^{m\times \sqrt{\chi \cdot m}\times m}$, sampled from
$\mu$, which together merge the two sketch legs of size $m$ into a
single one \cite[App.~C.1]{masolomonik2022}. We note that each tensor $\mathB_{i+1}$  is implicitly reshaped in the algorithm from $\R^{\chi\times m\times \chi}$  to $\R^{\chi \times \sqrt{m/\chi}\times\sqrt{\chi\cdot m}}\times \chi$. We perform our contraction matching \cite[App.~(C.1)]{masolomonik2022}, and each copy contracts to the sketched vector $\vv_{r}\in\R^{m}$.



\begin{algorithm}[H]
\caption{Sketched estimation of $\trace(\matA^{k})$ (MPS)}
\label{alg:mps}
\begin{algorithmic}[1]
\STATE \textbf{Input:} MPS cores $\tenB_i$, $i\in[N]$; power $k$; sketch size $m$.
\STATE \textbf{Output:} estimate $E_k$ of $\trace(\matA^{k})$.
\STATE \emph{Sampling convention:} for each copy $r\in[k]$, the Kronecker tensors $\ten{K}_n^{(r)}\in\R^{d\times m}$, $n\in[N]$, and the gadget pairs 
$g_{i}^{(r,a)}\in\R^{m\times\sqrt{m/\chi}\times m}$ and
$g_{i}^{(r,b)}\in\R^{m\times \sqrt{\chi \cdot m}\times m}$, $i\in[N-1]$ are drawn i.i.d.\ from $\mu$, Gaussian or Rademacher with mean $0$ and standard deviation $1/\sqrt{m}$.
\FOR{$r=1,\dots,k$}
  \STATE \emph{Kronecker step:} $\hat{\tenB}_{n}\gets\tenB_{n}\cdot\ten{K}_{n}^{(r)}$ for $n\in[N]$
  \STATE $\ten{W}\gets\hat{\tenB}_{1}$
  \FOR{$i=1,\dots,N-1$}
    \STATE $\ten{W}\gets\bigl(\bigl(\ten{W}\cdot g_{i}^{(r,a)}\bigr)\cdot\hat{\tenB}_{i+1}\bigr)\cdot g_{i}^{(r,b)}$
  \ENDFOR
  \STATE $\vv_{r}\gets\ten{W}\in\R^{m}$
\ENDFOR
\STATE $E_k\gets\prod_{r=1}^{k}\inner{\vv_r}{\vv_r}$
\RETURN $E_k$
\end{algorithmic}
\end{algorithm}

\begin{lemma}[Correctness and cost of \Cref{alg:mps}]\label{lem:mps-correct}
Given an MPS \eqref{eq:mps-def} with parameters $N,d,\chi$ representing
$\vb\in\R^{d^{N}}$, with $\matA=\vb\vb^{\top}$ and $k\ge1$ the output of
\Cref{alg:mps} is $E_{k}=\prod_{r=1}^{k}\inner{\vv_r}{\vv_r}$ (see Eq. \eqref{eq:estimator}) and the executed contraction tree is constrained on the left-to-right path $T_0$ (\Cref{def:constrained}). \\
For $m$ as in \Cref{cor:k1-gaussian} ($k=1$) or
\Cref{cor:k-gaussian} ($k\ge2$) with $d_{\Pi}=d^{N}$, it satisfies the
$(\varepsilon,\delta)$ guarantee Eq. \eqref{eq:goal}; and its total cost is
$O\!\bigl(kN(\chi m^{2.5}+\chi^{2}dm)\bigr)$, i.e. polynomial in $d,\chi,N,k$.
\end{lemma}
\begin{proof}
As explained before, from a theory perspective we look at the this algorithm as $
\bigl(\matS^{(r+1)}\vb\bigr)^{\top}
  =\matS^{(r)}\matA\,\matS^{(r+1)\top}$. Now the sketch $\matS^{(r)}=\mathrm{mat}(G_S^{(r)})$ assembled by \Cref{alg:mps} satisfies all the requirements for our  hypothesis of both \Cref{thm:k1} and \Cref{thm:tracek} as established in \cite[Thm.~3.1, App.~C.1]{masolomonik2022}.
So the output of \Cref{alg:mps} is exactly the estimator from Eq. 
\eqref{eq:estimator}, and as before, \Cref{cor:k1-gaussian} gives the $(\varepsilon,\delta)$ guarantee
of Eq. \eqref{eq:goal} for $k=1$ whenever
\[
  m=\Omega\!\left(\frac{N\bigl(N\log d+\log(1/\delta)\bigr)}
    {\varepsilon^{2}}\right),
\]
and \Cref{cor:k-gaussian} gives it for $k\ge2$ whenever
\[
  m=\Omega\!\left(\frac{N\,k^{2}\max\{1,\rho^{2}\varepsilon\}}
    {\varepsilon^{2}}\bigl(N\log d+\log(k/\delta)\bigr)\right).
\]
For the rank-one input $\matA=\vb\vb^{\top}$ each transfer block collapses to the
outer product $\matM_r=\vv_r\vv_{r+1}^{\top}$ with $\vv_r=\matS^{(r)}\vb$, so the
estimator is exactly $E_k=\prod_{r=1}^{k}\inner{\vv_r}{\vv_r}$.\\
Finally, the sketching contraction tree produced by \cite[Alg.~1]{masolomonik2022} is, by construction, constrained on the given data contraction tree.
\leanref{TensorTreeSketching.mps_estimator_eq_prod_normSq}\leanok
\leanuses{transferBlock_rankOne, estimator_eq_transfer_prod}

\medskip\noindent\emph{Computational Cost.} 
For each copy, the Kronecker step contracts $\ten{K}_n^{(r)}$ with the
physical leg of $\tenB_n$ at cost $\chi^{2}dm$ per site
(\Cref{rem:contraction-cost}), leaving the sketched nodes
$\hat{\tenB}_{n}$. Repeated $N-1$ times, contracting its 4 tensors in the prescribed order, at the cost.
Each iteration therefore costs $\chi m^{2.5}+\chi m^2$ per iteration (see  \cite[Eq.~C.3]{masolomonik2022} for further detail). 
We repeat it $k$ times and finish with $k$ dot products at cost $km$. The final cost is 
$$O(kN(\chi m^{2.5}+\chi^2dm)$$
\end{proof}
%--------------------------------------------------------------
% Experiments
%--------------------------------------------------------------
\section{Experiments}
\label{sec:experiments}

In this section we report experimental results with an implementation of our proposed algorithm. The goal of the experiments is to demonstrate the ability of .....

\subsection{Experiment 1}
Content goes here.

\subsection{Experiment 2}
Content goes here.


%--------------------------------------------------------------
% Conclusions
%--------------------------------------------------------------
\section{Conclusions}
\emph{[To be written.]} We gave a general framework --- theory and algorithms ---
for approximating $\trace(\matA^{k})$ of a PSD matrix $\matA$ in
tensor-network form. The central $(\varepsilon,\delta)$ guarantee is agnostic to the
structure of $\matA$, with sample complexity polynomial in $N$ and $k$; as
an explicit instantiation we treated the MPO, giving two sketches with efficient
contraction paths.


\section{Acknowledgments}
This research was supported by ...

%--------------------------------------------------------------
% Bibliography
%--------------------------------------------------------------

\printbibliography

%--------------------------------------------------------------
% Appendix A
%--------------------------------------------------------------
\appendix \section{Appendix Name}  
\label{app:label for appendix} 
text

\end{document}
